\begin{center}
\begin{tikzpicture}[x=1cm,y=1cm,font=\small,>=Stealth,
  every node/.style={inner sep=3pt},line cap=round,line join=round]
\path[use as bounding box] (-1,-1.6) rectangle (11,5.5);
\node[font=\small\bfseries,text=CStatement] at (5,5.1) {水平公共底：两段竖直高相加};

\coordinate (O) at (2,0); \coordinate (A) at (5,4);
\coordinate (B) at (3,2); \coordinate (E) at (3.5,2);
\fill[CStatement!18] (O)--(B)--(E)--cycle;
\fill[CExample!20] (A)--(B)--(E)--cycle;
\draw[thick] (O)--(A)--(B)--cycle;
\draw[very thick,CTrap] (B)--(E);
\draw[dashed] (2,0)--(6,0) (3.5,2)--(6,2) (5,4)--(6,4);
\draw[<->,CProof] (5.7,0)--node[right] {$2$}(5.7,2);
\draw[<->,CProof] (5.7,2)--node[right] {$2$}(5.7,4);
\node[below left] at (O) {$O(0,0)$};
\node[above] at (A) {$A(3,4)$};
\node[left] at (B) {$B(1,2)$};
\node[below right] at (E) {$E_h(\frac32,2)$};
\node[align=left] at (8.5,2.4) {$BE_h=\frac12$\\[6pt]竖直高：$2+2=4$\\[6pt]$S=\frac12\cdot\frac12\cdot4=1$};

\node[align=center] at (5,-1.15) {与竖直法使用同一三角形；公共底与高的方向互换。};
\end{tikzpicture}
\end{center}
